% 来源: 19c9fbb8-4ef2-8b58-8000-0000a991bfcb_2.8_有效传输_复用.pdf
% 提取时间: 2026-02-28T22:51:34+08:00

2.8 有效传输：复用
1. 复用技术概述
复用技术是数据通信中的基础方法，它使得多个信号能够共享单一传输介质，从而优化资源利用率。该
技术广泛应用于电信、计算机网络和广播系统中。




复用技术的主要目标
 ● 提高信道效率：允许多个用户共享带宽资源
 ● 降低基础设施成本：减少专用链路的需求
 ● 增强可扩展性：支持高需求网络(如蜂窝系统)的扩容



2. 复用技术分类
复用技术主要分为以下四大类型，根据信号分割维度的不同进行区分：
(1) 频分复用(FDM)
工作原理
将总带宽划分为多个互不重叠的子频带，每个信号独占一个频段传输。关键技术包括调制解调和带通滤
波。




                                                1
频域




时域
技术特性
 ● 采用模拟调制技术
 ● 需要设置保护频带防止干扰
 ● 频带分配固定且静态


典型应用
 ● 传统广播系统(AM/FM收音机)
 ● 有线电视网络(CATV)
 ● 早期电话系统(载波通信)

目前的AT&T与ITU FDM载波分配标准
 ● 基群（Group）：12 路语音信道 (每路4kHz) = 48kHz，范围从60kHz到108kHz
 ● 超群（Supergroup）：60路信道，420kHz至612 kHz之间载波上5个基群信号的FDM
 ● 主群（Mastergroup）：10个超群




                                                         2
OFDM
LTE采用OFDM主要有这么几点原因：
1、硬件技术的成熟，尤其是数字信号处理芯片技术的成熟，使OFDM实现成本和难度大幅降低。
2、OFDM更容易和MIMO结合。LTE提升速率主要依靠两点：第一是提升占用带宽（从5MHz到
20MHz），第二就是MIMO。
3、OFDM技术更干净，没高通CDMA专利限制。

迈向高数据速率时代
当数据传输速度很快时，由于信号通过不同路径到达接收端的时间有快有慢，前面的比特可能会和后面
几十个比特搅在一起，造成干扰。这种现象叫“符号间干扰”（ISI）。在很多情况下，这种干扰甚至可能
拖累几百甚至几千个比特。简单来说，那些在低速通信时没啥问题的信道，在高速通信时就变得特别糟
糕了！所以，我们必须想办法解决这个问题。
解决方案 - 均衡器(Solution - Equalizer)
贝尔实验室的 Robert Lucky 于 1964 年提出了一种有效解决此类问题的技术方案，即自适应均衡器
（Adaptive Equalizer）。作为一种滤波器，均衡器的主要功能是缓解信道衰落对接收信号（Rx 信号）
的影响，并消除符号间干扰。该系统以失真波形（Distorted Waveform）作为输入，该波形是原始信号

                                                       3
及其多径分量叠加的结果，经过处理后输出为清晰的目标比特流（Clean Desired Bit Stream），其工作
原理如下图所示。




           均衡器的输入是一个失真的波形，而其输出是一个干净的比特流

别把均衡器当成一个物理设备，均衡器其实就是微处理器里的几行代码。均衡器有它的优缺点。它大大
降低了误码率，让系统正常工作。但另一方面，在高速通信系统中，它是传统无线接收端里最耗资源、
要求最高的部分。具体消耗多少资源取决于系统的设计和需求，但它甚至可能占用 75% 的计算资源！在
高速通信中，如果你在一个时间窗口内忙着处理一个比特，而接下来的很多比特已经在接收端排队等着
处理，那问题就大了。这会导致缓冲区被填满的速度比清空的速度还快，显然是不现实的。
所以我们可以得出结论：低速通信对处理器的要求比较简单，而高速通信则需要性能更强的接收端处理
器。但这种高性能处理会消耗大量能量，对于像智能手机这样的电池供电设备来说并不实际。那么有没
有一种技术，可以用更简单的均衡器实现高速通信呢？答案是肯定的，那就是 OFDM。
时域中的 OFDM
正交频分复用(OFDM)也称为多载波调制，使用不同频率的多个载波信号,并在每个信道上发送一些比特。
在时域中，正交频分复用（OFDM）技术将一个高速串行比特流分解为多个低速并行比特流。随后，这
些并行的低速比特流分别调制到一组相互正交的正弦载波上。假设我们有个R bps的数据流，并且可用带
宽为 N fb 中心频率在 f0 。可以使用整个带宽来发送这个数据流，即每个比特的持续时间为1/R。另一
种方法是将该数据流分裂成N个子数据流，并使用串并转换器。每个子数据流的数据率为R/N bps，并通
     ​        ​




过一个独立的副载波来传送，相邻副载波之间相距 fb 。现在每个比特的持续时间就是N/R。其中，正
弦波之间的正交性定义为在特定时间间隔内其内积的积分或求和结果为零。这种正交特性使得各子载波
                             ​




能够在频谱上紧密排列而互不干扰，从而显著提高了频谱利用率。
通过示例详细阐述了这一过程：针对比特流 ( b[k] )，下图展示了以下操作步骤：



                                                         4
                                OFDM 示例
1. 比特流分解
   假设每个比特的持续时间为T，并需要传输总计 9 个比特。在高速通信系统中，发送所有比特所需
   的时间为 9T 秒。为了提高效率，我们将原始比特流分解为 9 条并行比特流，每条比特流的持续时
   间扩展为 9T 秒。
2. 正交子载波生成

   设基频为 f0 = 9T1 ，在此时间间隔内包含 9 个采样点（即一个完整周期）。随后，生成一组频率
            ​       ​




   分别为 f0 , 2f0 , 3f0 , … , 8f0 的正弦波信号。这些正弦波彼此正交，并构成如图所示的 9 个子载
   波（subcarriers）。
        ​       ​       ​   ​




3. 幅度调制
   接下来，将每个子载波与对应的比特值相乘，这种调制方式实现了比特信息到子载波上的映射。
4. 最后，将所有这些幅度缩放的正弦波相加，生成所需信号。请注意，此复合信号的持续时间为 9 T
   秒（与原始比特流的持续时间相同），但包含所有信息。

一、OFDM时域信号的数学表达

                                                              5
假设OFDM系统包含 ( N ) 个子载波，每个子载波的中心频率为
 fk = f0 + k ⋅ Δf , k = 0, 1, … , N − 1 ， Δf 为子载波间隔。
对于第 k 个子载波，其携带的调制符号为 dk (如QAM或PSK调制后的符号），则单个OFDM符号的时
     ​       ​




                                                         ​




                     N −1
域表达式为： s(t) = ∑ dk ⋅ g(t) ⋅ ej2πf t , t ∈ [0, T ]
                             ​   ​
                                                 k   ​




                     k=0
其中：
  ● T 为OFDM符号的持续时间；
  ● g(t) 为矩形窗函数（理想情况下），即 g(t) = 1 （ t ∈ [0, T ] ），否则为升余弦窗等成形函
    数；
  ● ej2πf t 为第 k 个子载波的复指数载波信号。
                 k



二、子载波正交性的时域本质
OFDM的关键特性是子载波在时域上正交，即任意两个不同子载波 fk 和 fm k = m 符号周期
T 内的内积为零：
                                                                 ​   ​




∫ T ej2πfk t ⋅ e−j2πfm t dt = ∫ T ej2π(k−m)Δf t dt = 0 (k =
         ​
                     ​   ​




                                             ​             m)
 0                                       0
正交性的时域直观理解：
 ● 当子载波间隔 Δf =         时，任意两个不同子载波在符号周期 T 内恰好相差整数个周期，因此它
                    1
                                     ​




                    T
   们的乘积在 T 内的积分（能量）为零
 ● 这一特性使得接收端可以通过匹配滤波或**离散傅里叶变换（DFT）分离各子载波信号，避免ICI。

图解
对限制在 [0, 2π] 内的 sin(t) 信号，形成过程如下




                                                                         6
对限制在 [0, 2π] 内的 sin(2t) 信号，形成过程如下




将 sin(t) 和 sin(2t) 信号频谱叠加在一起，形成OFDM的频谱




OFDM有几个优点。首先，频率选择衰减仅影响一些副载波，而不是整个信号。如果据流被前向纠错码
保护，此类衰减就很容易解决。更重要的是，OFDM克服了多径环境中的码间干扰(ISI)。比特率越高，码
间干扰的影响也就越大，因为比特间或符号间的距离越小。使用OFDM，数据率会降低1/N倍，也就是说
符号周期增加了N倍。因此，如果源数据流的符号周期是 Ts ，OFDM信号的周期就是 N Ts 。这就急
剧降低了码间干扰的影响。作为一项设计技巧，我们要选择 N ，使 N Ts 远远大于信道带来延迟的均
                                    ​           ​




方根值。根据上述考虑，在使用OFDM后，也许就没有必要再应用均衡器了，因为它是一种很复杂的设
                                         ​




备，并且带来码间干扰的符号的数量越多，它就越复杂。
与OFDM一起使用的常见调制机制是四相相移键控(QPSK)。在这种情况下，每个被传送的符号用2比特
来表示。OFDM QPSK机制的一个例子是由512个载波构成的6 MHz带宽，载波距略小于12 kHZ。为了使
码间干扰最小化，数据以突发形式传送，每次发送由一个循环前缀及其后紧跟的数据符号构成。这个循

                                                      7
环前缀用于吸收因多径而引起的上一次发送的不稳定性。对于这个系统，每次突发包括由64个符号构成
的循环前级，以及之后的512个QPSK符号。因此在每个副载波上，QPSK符号被持续时间为64/512个符
号时间的前缀分隔开，这个前缀又称为循环前缀(cyclic prefix，CP)。一般来说，当前缀发送完毕时，由
组合的多径信号产生的波形不会是前一次发送的任何样本的函数，因此就不存在码间干扰。
OFDM信号处理涉及两个函数，称为快速博里叶变换(FFT)和快速傅里叶逆变换(IFFT)。FFT是将一组均匀
分布的数据点从时域变换到频域的算法。实际上，FFT是能够进行快速数字傅里叶变换的一个算法族，它
们构成了离散傅里叶变换(DFT)中的一种特殊情况。DFT是指任何能够生成时域函数的量化傅里叶变换的
算法。IFFT是FFT的逆运算。对OFDM来说，源比特流被映射到一组M个副载波频带上。然后，为了生成
发送信号，每个副载波都要执行IFFT操作以生成M个时域信号。这些信号被依次矢量相加，并生成用于
发送的最终时域波形。IFFT操作具有保证副载波互不干扰的效果在接收端，使用FFT模块将输人信号映射
回M个副载波，从中就可以恢复这些数据流了。
(2) 时分复用(TDM)
基本概念
将时间轴划分为固定长度的时隙帧，每个连接周期性占用特定时隙。可分为同步TDM和统计TDM两种形
式。




同步TDM
 ● 固定时隙分配：数据被组织成“帧”：每帧包含一组循环使用的时隙；每个数据源可以被分配一个或
   多个时隙
 ● 严格的时钟同步要求：同步时分复用中同步是指时隙被提前分配给数据源且是固定的。
 ● 存在空闲时隙浪费：如果没有信息，输入设备将发送空时隙。复用的线路上的数据率是固定的。
 ● 同步时分复用 Synchronous time division multiplexing：可以用于数字信号或模拟信号传输数字数
   据。
 ● 脉冲填充：TDM最困难的问题是同步不同的数据源，输入数据率之间也可能不存在简单的比例关
   系，需要在每个输入信号中填充额外的比特或脉冲，直到输入速度被提高到本体时钟信号速率，如
                                                                  8
  下所示




典型应用场景
ITU对于TDM多路电话通信系统，制定了两种准同步数字体系(PDH)和两种同步数字体系(SDH)标准的建
议。
准同步数字体系(PDH)，ITU提出的两个建议：
  ● E体系 PCM30/32路数字载波系统，以2.048Mbps作为基群速率(E1速率) 的数字系列 。欧洲、中国
    采用，并用于国际间传输
  ● T体系 PCM24路数字载波系统，以1.544Mbps作为基群速率(T1速率) 的数字系列。用于北美、日本

我们以E体系为例简单介绍。该体系以2.048Mbps的E1电路为基础，采用时分复用（TDM）技术实
现多路语音和数据信号的数字化传输。E体系采用分层复用结构，每一级将多个低速率信号复合成
更高速率的信号。具体包括：E1（2.048Mbps，30路64Kbps语音信道）、E2（8.448Mbps，4
个E1）、E3（34.368Mbps，4个E2）、E4（139.264Mbps，4个E3）和E5（565.148Mbps，4
个E4）。系统采用PCM编码和CRC校验技术，确保传输质量。E体系具有严格的同步要求，采用
字节交织复用方式，广泛应用于传统电信网络、SDH传输系统等领域。

                                                             9
E体系一次群（E1）帧结构详解




1.基本参数
● 帧长：256 bits（32时隙×8 bits）
● 帧频：8000帧/秒（125μs/帧）


                             10
 ● 总速率：32×8×8000 = 2.048 Mbps
● 有效话路：30路（时隙1-15，17-31）

2. 时隙分配

时隙号                      用途                说明
TS0                      帧同步               交替传送帧同步码（0011011）
                                           和对端告警信息
                                           E1的TS0同步码在偶数帧为
                                           0011011，奇数帧保留第2比特
                                           传告警信息，其他比特固定为1
TS1-TS15                语音信道（第1-15路）       每时隙64Kbps PCM编码
TS16                    信令时隙               采用随路信令（CAS）方式，
                                           通过复帧传送30路的信令
TS17-TS31               语音信道（第16-30路）      每时隙64Kbps PCM编码
 3.复帧结构
 ● 1个复帧 = 16个基本帧（F0-F15）
 ● 信令传输：
    ○ F0的TS16传送复帧同步码（0000）
    ○ F1-F15的TS16传送30个话路的ABCD信令位（每个话路占用4bit）


PDH向SDH/SONET的演进过程
1. PDH时代的局限性（1970s-1980s）
PDH（准同步数字体系）是早期数字传输的核心技术，采用异步复用（各节点时钟独立，允许±50ppm
偏差）和逐级解复用机制，存在以下根本性缺陷：
  ● 制式割裂：欧洲（E1，2.048Mbps）、北美（T1，1.544Mbps）、日本（J1，1.544Mbps）互不兼
    容，跨国通信需复杂转换。
  ● 运维低效：从140Mbps信号中提取2Mbps需层层解复用（如E4→E3→E2→E1），故障定位困难。
  ● 管理薄弱：仅2%的开销字节（如E1的TS0用于同步），无法实现端到端性能监控。

2. SDH/SONET的标准化（1980s末）
为突破PDH限制，国际电信联盟（ITU-T）和ANSI分别主导制定：
  ● SDH（1988年）：基于ITU-T G.707标准，以**STM-1（155.52Mbps）**为基本模块，统一欧洲/


                                                               11
    亚洲市场。
  ● SONET（1984年）：基于ANSI T1.105标准，以**STS-1（51.84Mbps）**为基础，主导北美市
    场。
3. 关键技术革新
SDH/SONET通过以下创新解决PDH痛点：
 技术维度 PDH                      SDH/SONET
 同步机制 准同步（各节点独立时 全网同步（原子钟+SSU，精度达10⁻¹¹）
            钟）
 复用方式 异步比特填充                   同步字节间插（直接分插低阶信号）
 业务映射 硬性速率适配                   **虚容器（VC）**灵活封装PDH业务（如VC-12对应
                               2Mbps）
 管理开销 2%帧开销                    5%段/通道开销（B1/B2/B3字节实现性能监测）
 保护倒换 无保护或人工切换                 自愈环（MS-SPRING，倒换时间<50ms）
4. SDH与SONET的对应关系
 SDH等级 SONET等级 速率（Gbps）        光波长（nm）      典型应用场景
 STM-1 OC-3       0.155        1310/1550    城域网接入、基站回传
 STM-4 OC-12      0.622        1310/1550    省内骨干网、数据中心互联
 STM-16 OC-48     2.488        1550         国家骨干网、海底光缆
 STM-64 OC-192    9.953        1550（DWDM）   超高速骨干网、100G以太网承载
 STM-256 OC-768   39.813       1550         未来400G/1T网络
                               （Coherent）
 5.总结
SDH（STM-N）和SONET（OC-N）是同步数字传输的两大标准体系，通过分层复用实现从155M到
40G的速率覆盖。二者在高速等级（≥155M）可互通，但基础模块和映射方式存在差异。当前
SDH/SONET仍是传统专网的重要技术，但正逐步向更高效的OTN和分组传输演进。
统计TDM

                                                                  12
 ● 动态时隙分配
 ● 支持可变长度数据单元
 ● 更高的带宽利用率

统计时分复用 Statistical Time Division Multiplexing




典型应用场景
统计时分复用（Statistical Time Division Multiplexing，StatTDM）是一种动态分配时隙的复用技术，
适用于突发性数据流量场景。其核心特点是按需分配带宽，相比固定时隙的同步TDM（如E1/T1），能显
著提高信道利用率。以下是典型应用场景：
1. 计算机网络
以太网（Ethernet）：采用统计复用共享总线/交换机端口，动态分配带宽给不同主机，支持
10M/100M/1G等速率自适应。
IP网络：路由器通过统计复用处理不同来源的IP包，实现互联网“尽力而为”传输。
2. 无线通信
4G/5G空口资源调度：基站动态分配时隙和频段给用户（如5G NR的Slot调度），适应移动业务流量的
突发性。
Wi-Fi（802.11协议）：CSMA/CA机制本质是统计复用，允许多终端竞争信道。
3. 数据中心与云计算
                                                                   13
虚拟化网络：虚拟机（VM）间通过vSwitch统计复用物理网卡带宽（如AWS Nitro系统）。
SDN流量工程：控制器动态调整路径带宽分配，应对流量潮汐效应。
4. 语音/视频传输
VoIP（如SIP协议）：将语音包统计复用到IP网络，静默期可释放带宽给其他业务。
视频会议系统：H.323/Skype等动态分配带宽，适应视频码率波动。
同步TDM vs StatTDM 对比
特性                     统计TDM（StatTDM）        同步TDM（如E1）
时隙分配                   动态（按需）                固定（预分配）
适用流量                   突发性数据（IP包）            恒定速率业务（语音）
信道利用率                  高（可超售带宽）              低（存在空闲时隙）
典型协议                   以太网、IP、MPLS           PDH/SDH、T1/E1

(3) 波分复用(WDM)
基本原理
在光纤通信中利用不同波长的光载波传输多路信号，本质上是光域的频分复用技术。




将一束光分解成它的各种颜色           使用耦合器和滤波器的WDM             光复用/解复用器
密集波分复用( DWDM)
密集波分复用（DWDM，Dense Wavelength Division Multiplexing）是一种先进的光通信技术，通过在
单根光纤中同时传输多个波长不同的光信号来大幅提升传输容量。它将光纤的可用光谱划分为多个紧密
间隔的波长通道（通常间隔0.8nm或更小），每个通道可承载独立的数据流，从而实现单光纤传输容量
                                                                14
从数十Gbps到数Tbps的飞跃。DWDM系统的核心组件包括可调谐激光器、光复用器/解复用器和掺铒光
纤放大器（EDFA），这些设备协同工作确保信号在长距离传输中的质量和稳定性。该技术广泛应用于长
途骨干网、海底光缆和数据中心互联等场景，是构建现代高速光网络的基础。
DWDM技术的主要优势在于其极高的频谱利用率和灵活的扩容能力。相比传统的CWDM（粗波分复用）
技术，DWDM的波长间隔更小，可在C波段（1530-1565nm）和L波段（1565-1625nm）内支持160个
以上的波长通道。系统采用数字相干检测技术，结合先进的调制格式（如QPSK、16-QAM），使单波长
速率达到400Gbps甚至更高。此外，DWDM支持灵活栅格（Flexi-Grid）技术，允许动态调整通道间隔
和带宽分配，以适应不同业务需求。随着光器件成本的降低和技术的成熟，DWDM正逐步向城域网和接
入网延伸，为5G、云计算和物联网等应用提供强有力的传输支撑。
 特性                 DWDM                CWDM
 波长间隔               0.4~0.8nm（密集）       20nm（稀疏）
 通道数                80~160（C+L波段）       16~18
 传输距离               1000km+（带放大器）       ≤80km（无放大器）
核心技术组件
 ● 光源：
    ○ 可调谐激光器（ITLA）：波长精确可控（±0.01nm）。
    ○ 调制格式：NRZ（10G）、DP-QPSK（100G）、16-QAM（400G+）。
 ● 光复用/解复用器：
    ○ 阵列波导光栅（AWG）：实现多波长合分波。
    ○ 薄膜滤波器（TFF）：低插损，用于通道选择。
 ● 光放大器：
    ○ 掺铒光纤放大器（EDFA）：C波段（1530~1565nm）放大。
    ○ 拉曼放大器：分布式增益，扩展至L波段（1565~1625nm）。

应用场景
 ● 长途干线：单纤容量可达48Tbps（160λ×300G）。
 ● 数据中心互联（DCI）：低成本简化版DWDM（如OpenZR+）。
 ● 5G前传/中传：通过DWDM承载eCPRI（25G/100G波长）。



(4) 码分复用(CDM，Code Division Multipling)

                                                       15
码分复用（Code Division Multiplexing, CDM）是一种通过正交编码实现多路信号共享同一频带
和时间的复用技术。其核心思想是：
 ● 每路信号分配唯一的扩频码（Spreading Code）
 ● 所有信号在相同频段上同时传输
 ● 接收端通过相关检测分离目标信号


核心技术
正交扩频码是一组具有良好互相关特性的码序列，在码分多址（CDMA）系统中区分不同用户的信
号。其核心特点是：
 ● 正交性：不同用户的扩频码在同步条件下互相关为零（或极低），从而减少多用户干扰
   （MAI）。
 ● 扩频增益：通过将窄带信号扩展到更宽的频带，增强抗干扰和抗多径能力。




通过一个例子来解释正交扩频码的作用。如上图，假设A发送了一个1，而其他信道处于离线状态且没有
噪声。那么接收方将接收到<1,-1,-1,1,-1,1>。这是通过对接收到的信号与A的编码进行直接乘积来解码
的：
<1,-1,-1,1,-1,1> * <1,-1,-1,1,-1,1>
= (1*1) + (-1 * -1) + (-1 * -1) + (1 * 1) + (-1 * -1) + (1 * 1) = 6
这里的6表明非常有信心A发送的是1。如果A发送的是-1，那么直接乘积的结果将是-6。
但是，其他信道呢？信道B的解码是通过计算接收到的信号与B信道的码片序列的直接乘积来完成的。这
给我们：
<1, -1, -1, 1, -1, 1> * <1, 1, -1, -1, 1, 1> = (1) + (-1) + (-1*-1) + (1*-1) + (-1) + (1) = 0


                                                                                           16
因此，为信道A发送的码型序列没有为信道B生成一个值。这是因为信道A的码片序列与信道B的码片序列
正交（线性无关）。

如何生成正交扩频码呢，Walsh码（Walsh-Hadamard Code）是一个典型的生成方法。
1. 基本原理
Walsh码是一种基于哈达玛矩阵（Hadamard Matrix）的正交码，其生成方式如下：
  ● 递归构造：
     ○ 1阶哈达玛矩阵： H1 = [1]
                      ​   ​




     ○ 2n 阶哈达玛矩阵：




                                            ,
    ○   递归实现，可以扩展正交扩频码
Walsh具有如下特点
 ● 正交性：Walsh码最重要的特性是正交性，即任意两个不同的Walsh码的内积为0。例如，对于
                                n−1
    Walsh码 Wi 和 Wj （ i = j )，有 ∑ Wi (k)Wj (k) = 0 。这种正交性使得Walsh码在多
           ​      ​             ​   ​   ​




                                k=0
    址通信中可以有效地分离不同用户的信号，减少多址干扰。
 ● 自相关性：Walsh码的自相关性具有良好的特性。对于一个Walsh码 (W)，其自相关函数在码元同步
    时为最大值（等于码长），在非同步时为0或很小的值。这有助于在接收端进行信号的同步和检测。
 ● 可分性：Walsh码可以按照不同的长度进行分组，较长的Walsh码可以由较短的Walsh码通过一定的
    规则生成。这种可分性使得Walsh码在不同的系统带宽和数据速率要求下具有较好的灵活性。
还有其他的扩频码
 特性               Walsh码            PN码（伪随机码）          Gold码
 正交性              严格正交（同步时） 非正交                        准正交
 长度               2^n               任意长度               2^n−1
 抗干扰              依赖同步              强（异步适用）            中等
 典型应用             CDMA下行链路          加密、抗截获             卫星导航（GPS）


                                                                 17
这是因为Walsh码存在以下一些不足，这使得在一些通信系统中需要结合PN码（伪随机噪声码）来使
用：
 ● 非理想的多址干扰抑制：虽然Walsh码在同步条件下具有良好的正交性，能有效区分不同用户信号，
   但在实际通信环境中，由于多径传播、多普勒频移等因素，信号到达接收端时往往存在不同步的情
   况。此时，Walsh码的正交性会遭到破坏，产生多址干扰，影响系统性能。
 ● 码资源有限：Walsh码的数量是有限的，其数量等于码长。例如，在IS - 95系统中使用64阶Walsh
   码，只能提供64个不同的码字。随着系统中用户数量的增加，码资源会逐渐紧张，当用户数超过码
   资源数量时，就无法为每个用户分配唯一的Walsh码，限制了系统容量。
 ● 对信道变化适应性差：Walsh码是固定的正交码组，一旦确定就难以根据信道的实时变化进行调整。
   而实际通信信道是时变的，信道条件会不断发生变化，如信号衰落、干扰变化等。Walsh码不能很好
   地适应这种变化，无法根据信道状态动态地优化信号传输，降低了系统的性能和效率。
相比之下，PN码具有良好的自相关性和类似噪声的特性，在不同步时也能提供较好的多址干扰抑制能
力，且可以通过不同的序列生成方式提供大量的码资源，还能根据系统需求灵活调整码序列，更好地适
应信道变化。因此，在CDMA等通信系统中，通常将Walsh码和PN码结合使用，以充分发挥它们各自的
优势，提高系统的性能和可靠性。例如，在CDMA系统中，Walsh码用于在前向链路中区分不同的信道，
而PN码用于扩频调制和在反向链路中区分不同的用户。
实现方式
扩频是无线通信的一种重要的编码形式。扩频既可用于传输模拟数据，也可用于传输数字数据。扩频技
术最初是针对军事及情报部门的需求而开发的。它的基本思想是将携带信息的号扩散到较宽的带宽中，
用以加大干扰及窃听的难度。开发出的第一种扩频类型称为跳频（跳频扩频技术是由Hedy Lamarr在
1940年发明的U.S. Patent 2292387，当时26岁）。较新一些的技术是直接序列扩频。这两种技术广泛
应用于各种无线通信标准和产品中。




上图重点描绘了所有扩频系统的主要特点。输人的数据进入信道编码器，以生成模拟信号，这个模拟信
号具有围绕某个中心频率的、相对较窄的带宽。然后这个信号被称为扩频码或扩频序列的数字序列进一
步调制。通常(但并非绝对)情况下，这个扩频码由伪随机样本或伪随机数生成器产生。此次调制带来的影


                                                        18
响是传输信号的带宽有显著增加(频谱被扩展)。在接收端，使用相同的数字序列对扩频信号进行解调。最
后，信号进入信道解码器，被还原成数据。
从这种明显的频谱浪费中可以得到以下几个好处。
 ● 我们可获得对各种噪声和多路失真的抗扰性。扩频最早应用在军事上，就是利用了它对人为干扰的
   抗干扰能力。
 ● 它也可用于信号的隐蔽和加密。只有知道扩频码的接收方才能恢复加密过的信息。
 ● 多个用户可以做到彼此之间很少干扰地独立使用相同的带宽。这个特点被用于蜂窝电话的应用中，
   它所使用的技术就是码分多址。
1）直接序列扩频(DSSS)
直接序列扩频（Direct Sequence Spread Spectrum，DSSS）是一种常用的扩频通信技术，以下从原
理、特点及应用等方面进行介绍：
基本原理
 ● 扩频过程：发送端将待传输的原始信号（通常是速率较低的数字信号）与一个高速伪随机序列（PN
   码）进行模二加（或相乘）运算，从而将原始信号的频谱扩展到一个很宽的频带上。例如，原始信
   号的频谱带宽为 (B_1)，PN码的速率为 (R_c)，且 (R_c\gg B_1)，经过扩频后信号的带宽将扩展到
   与 (R_c) 相关的较宽频带 (B_2)（(B_2\gg B_1)）。
 ● 解扩过程：在接收端，使用与发送端相同的PN码与接收到的扩频信号进行模二加（或相乘）运算，
   将扩频信号还原为原始的窄带信号。由于只有与发送端同步的正确PN码才能完成解扩，其他干扰信
   号在解扩后仍保持在宽带状态，通过后续的窄带滤波等处理，可以有效抑制干扰，恢复出原始信
   号。




                                                            19
2）跳频扩频(FHSS)
跳频扩频（Frequency - Hopping Spread Spectrum，FHSS）是另一种重要的扩频通信技术，以
下从原理、特点和应用等方面进行介绍：
基本原理
 ● 频点随机跳变：发送端在伪随机序列（PN 码）的控制下，按照一定的规律在一组预先指定的
   频率点上随机跳变地发射信号。例如，通信系统预先设定了一个包含多个频点的频率集合，
   PN 码会指示发射机在不同的时刻选择不同的频点进行信号发射，使得原始信号的频谱在这些
   频点之间快速跳变，从而将信号扩展到一个较宽的频带上。
 ● 接收端同步跟踪：接收端需要与发送端保持精确的同步，使用相同的 PN 码和跳频图案来跟踪
   发送端的频率跳变。在每个跳变时刻，接收端能够准确地将本地接收机的频率调谐到与发送端
   相同的频点上，以便正确接收信号。然后通过解跳和解调等处理过程，恢复出原始的信息信
   号。
通常采用多频移键控（MFSK）技术。该技术的特点是信号的频率每Tc秒变化一次，而每个信号元素的
持续时间为Ts秒。慢跳频扩频（Slow FHSS）系统的特征为Tc远大于Ts；相比之下，快跳频扩频（Fast
FHSS）则具有Tc小于Ts的特性。FHSS因其对噪声和干扰的高度抵抗性而著称，尤其是快速跳频系统，
在性能表现上更为优越。


                                                           20
Slow FHSS




Fast FHSS
(5) 空分复用(SDM)
空分复用（Space-Division Multiplexing，SDM）是一种用于提高通信系统容量和频谱效率的技术，以
下从原理、实现方式和应用等方面进行介绍：
基本原理
空分复用是利用空间位置的不同来区分不同的信号，在同一时间和频率资源上同时传输多个独立的信
号，从而提高系统的传输容量。其核心思想是将通信空间划分为多个相互独立的子空间，每个子空间可
以独立地传输不同的信息，使得多个用户或数据流能够在同一频段上同时进行传输而互不干扰。
实现方式

                                                           21
 ● 多天线系统：在发射端和接收端分别部署多个天线，形成多输入多输出（MIMO）系统。通过对不同
   天线发送的信号进行特定的编码和处理，使得接收端能够利用信号在空间传播中的不同特性，如到
   达角度、幅度、相位等，来区分不同天线发送的信号。例如，在一个2×2的MIMO系统中，发射端有
   两根天线，接收端也有两根天线，发射端可以同时通过两根天线发送两个不同的数据流，接收端利
   用信号的空间特征将这两个数据流分离并正确解调。
 ● 智能天线技术：在智能天线系统中，通过波束赋形可以将信号集中发送到用户所在方向，增强
   用户接收到的信号强度，减少对其他用户的干扰。采用具有方向性的天线阵列，通过自适应算法
   调整天线的波束方向，使其指向不同的用户或信号源。这样，不同方向的信号可以在空间上被分
   离，实现空分复用。例如，在移动通信基站中，智能天线可以根据手机用户的位置动态调整波束，
   将信号集中发送到用户所在方向，同时抑制其他方向的干扰信号，提高系统的容量和信号质量。
MIMO
MIMO技术通过提高频谱效率实现了更高数据速率的承诺。由于系统性能的潜在改进以及数字信号处理的
发展，许多无线系统如 IEEE 802.11n WLAN、基于 IEEE 802.16e 的 Mobile WiMAXTM Wave 2 和长
期演进 (LTE) 等移动无线系统近来已经采用了 MIMO技术和多天线技术。所有这些商用无线系统都工作
在高度多径环境下，正是这种多径环境保证了在使用多天线配置时的性能改善。
尽管在多径丰富的环境下运行时，MIMO 具有增强信号鲁棒性和提高容量的潜力，但是 MIMO器件与系
统的开发和测试需要一些高级信道仿真工具，这些工具应当易于配置，并可以准确地表示实际无线信道
和条件。
本文首先回顾 MIMO技术和无线信道的基本特性，然后介绍空间相关概念及其对 MIMO性能的影响。还
包括对 MIMO信道空间特性建模的示范，并描述如何使用商用仪表 (例如信号源) 对这些复杂信道进行仿
真。
从SISO到MIMO多路输入多路输出
通过充分利用无线信道的 "空间" 特性，可以使用布置在无线通信系统中发射机和/或接收机处的多根天
线，实质性地提高系统性能。这些系统，现在广泛称为 "多路输入多路输出 (MIMO)"，即在发射机和接
收机处设置两根或更多根天线。
在學習MIMO之前，需要先介绍一下什么是SISO, SIMO以及MISO。
SISO（Single-Input Single-Output）
在介绍MIMO之前，需要先介绍一下什么是SISO。从字面上理解，SISO 就是单发单收，是一种单输入单输出系
统，发射天线和接收天线之间的路径是唯一的，传输的是1路信号。在无线系统中，我们把每路信号定义为1个空
间流（Spatial Stream）。


                                                                     22
SISO示意图
由于发射天线和接收天线之间的路径是唯一的，这样的传输系统是不可靠的，而且传输速率也会受到限制。
SIMO（Single-Input Multiple-Output）
为了改变这一局面，在终端处增加1个天线，使得接收端可以同时接收到2路信号，也就是单发多收。这样的传输
系统就是单输入多输出，即SIMO。




SIMO示意图
虽然有2路信号，但是这2路信号是从同一个发射天线发出的，所以发送的数据是相同的，传输的仍然只有1路信
号。这样，当某一路信号有部分丢失也没关系，只要终端能从另一路信号中收到完整数据即可。虽然最大容量还
是1条路径，但是可靠性却提高了1倍。这种方式叫作接收分集。
MISO（Multiple-Input Single-Output）
我们换一个思路，如果把发射天线增加到2个，接收天线还是维持1个，会有什么样的结果呢？




MISO示意图
因为接收天线只有1个，所以这两路最终还是要合成1路，这就导致发射天线只能发送相同的数据，传输的还是只
有1路信号。这样做其实可以达到和SIMO相同的效果，这种传输系统叫作多输入单输出，即MISO。这种方式也
叫发射分集。
MIMO
如果收发天线同时增加为2个，那么是不是就可以实现独立发送2路信号、速率翻倍了呢？答案是肯定的，因为
从前文对SIMO和MISO的分析来看，传输容量取决于收、发双方的天线个数。而这种多收多发的传输系统就是
MIMO。

                                                   23
MIMO示意图
MIMO 技术允许多个天线同时发送和接收多个信号，并能够区分发往或来自不同空间方位的信号。通过空分复用
和空间分集等技术，在不增加占用带宽的情况下，提高系统容量、覆盖范围和信噪比。
MIMO技术的优点
MIMO技术的优点是通过增大天线的数量来传输信息子流，将多个数据子流同时发送到信道上，各发射信
号占用同一频带，从而在不增加频带宽度的情况下增加频谱利用率。使用MIMO技术后，可以令无线信号
的传输距离、天线的接受范围进一步扩大，信号抗干扰性更强，无线传输更为精准快速。而前面针对提
升手机在LTE或WiFi环境下的通讯品质，而推出的新式MIMO架构天线方案，无疑为用户带来了更流畅、
快速的无线应用体验。
MIMO技术通过空分复用和空间分集等技术，在不增加占用带宽的情况下，提高系统容量、覆盖范围和信
噪比。MIMO与传统的单天线系统相比多个发射和接收天线为无线系统的设计者打开了一个新的维度--
空间自由度。信号在多对收发天线间经历不同的信道衰落，如果这些衰落的统计特性互相独立，就相当
于在通信系统中引入了多个传输通道。这和增加系统传输带宽几乎可以达到同样的效果。上世纪90年代
贝尔实验室一篇介绍 ‘MIMO V-BLAST’技术的论文引发了学术界MIMO技术研究的热潮。在2000 年至
2010 年，MIMO 技术逐步在移动通信网络得到应用，成为3G 和4G 的重要组成部分。在3G 网络中，
MIMO技术通过分集技术增强传输链路的抗衰落能力，为早期蜂窝通信的稳定性提供了技术支持。随着
4G 长期演进技术（Long Term Evolution，LTE）标准的引入，多用户MIMO 技术实现了对多个用户的
并行传输，显著提升了频谱效率和系统容量。在这一阶段，MIMO 技术被广泛应用于高清视频传输和高
速数据业务场景，成为支撑高速传输的重要技术手段。进入2010 年，大规模MIMO技术应运而生，其核
心思想是在基站端部署数百甚至上千根天线，实现了小尺度衰落的信道硬化和渐进正交，并通过低复杂
度预编码技术，有效减少了用户间的互干扰，显著提升了系统的频谱效率和能量效率。在5G 通信网络
中，大规模MIMO 已成为增强型移动宽带、超
可靠低时延通信和大规模物联网连接场景的核心技术。例如，通过波束赋形技术，大规模MIMO 能够在
毫米波频段有效对抗路径损耗，为工业物联网等高可靠性场景提供毫秒级时延和稳定的连接。自2020 年
以来，MIMO 技术进一步向高维扩展和智能化方向演进。超大规模MIMO（XL-MIMO）在大规模MIMO
的基础上扩展天线阵列尺寸，提高了空间分辨率，大幅提升了空间复用和空间分集增益，可支持更高密
度的用户接入和更加精确的波束赋形，适用于超密集连接和高精度定位场景，是未来6G 高谱效传输的关
键技术之一。
Massive MIMO

                                                         24
Massive MIMO（大规模天线技术，亦称为Large Scale MIMO）是第五代移动通信（5G）中提高系统容
量和频谱利用率的关键技术。它最早由美国贝尔实验室研究人员提出，研究发现，当小区的基站天线数
目趋于无穷大时，加性高斯白噪声和瑞利衰落等负面影响全都可以忽略不计，数据传输速率能得到极大
提高。对于Massive MIMO，我們可以从两方面理解：
（1）天线数 - 传统的TDD网络的天线基本是2天线、4天线或8天线，而Massive MIMO指的是通道数达
到64/128/256个。
（2）信号覆盖的维度 - 传统的MIMO我们称之为2D-MIMO，以8天线为例，实际信号在做覆盖时，只
能在水平方向移动，垂直方向是不动的，信号类似一个平面发射出去，而Massive MIMO，是信号水平维
度空间基础上引入垂直维度的空域进行利用，信号的辐射状是个电磁波束。所以Massive MIMO也称为
3D-MIMO。

4. 总结
复用技术作为通信系统的核心支撑，其发展历程体现了从简单频分到多维复用的演进路径。现代通信系
统往往采用混合复用策略，如5G网络同时整合TDD、OFDMA和SDMA技术。理解各类复用技术的原理和
适用场景，对通信系统设计和优化具有重要意义。




                                                         25
